\scanpage{49}
\mspage{47}

que si on \unsure{un|a un} hom. de ind-objets \uline{de Cat}, \unsure{ou.|ou,} ce qui revient maintenant au même, un hom \deleted{de} Cat

\[
K_\infty \xrightarrow{u_\infty} K'_\infty
\]

tel que $\forall n,\quad u_\infty(K_n)\subset K'_{\deleted{\unsure{\tau(n)}}n'}$ pour $n'=\tau(n)$ convenable, \deleted{et si \unsure{cet hom.|cet} \unsure{est tel|est} que} et si on suppose que l'hom de ind-objets induit un isom. pour les foncteurs ind-représentés \uline{sur Cat hot} (p. ex. si $\underline{\underline K}$ et $\underline{\underline {K'}}$ \inserted{ind-}représentent le $=$ foncteur $H$, et l'hom de ind-objets est celui qui correspond à $\mathrm{id}_H$), alors $u_\infty$ est forcément un isom. \uline{dans Hot}, donc \unsure{aussi} \inserted{(sauf erreur)} dans Cat hot puisque $K_\infty,\ K'_\infty$ sont des catégories de Thomason-Kan. Mais ce qui n'est pas évident, c'est l'existence d'un hom $u_\infty$ comme dessus, et s'il existe, qu'il soit unique \inserted{(ou du moins canonique ?)} comme flèche de Cat hot.

Pour l'existence d'un $\unsure{\dot u_\infty|u_\infty}$, on est ramené à prouver que si on a un diagramme à flèches pleines

\diagram{\begin{tikzcd}[ampersand replacement=\&,column sep=large]
K_n \arrow[r,"i_n","{\mathrm{cofib.}}"'] \arrow[d,"u_n"] \& K_{n+1} \arrow[dr,dashed,"u_{n+1}"] \& {}\\
K'_{\tau(n)} \arrow[rr,"{\text{transition}}"' ,"{\unsure{\scriptscriptstyle\blacktriangleright|\cdot}}" marking] \& {} \& K'_N
\end{tikzcd}},

et si une flèche $u_{n+1}$ est donnée \inserted{``en pointillés''} \uline{dans Cat hot}, de façon à rendre le carré commutatif,

% archive sheet number: 48
